%% kapa: exam=1 round=36 year=2025 period=2 form=A subject=appraisal_law subject_ko=감정평가관계법규 seq=8 number=8
%% source: Q-Net 감정평가사 기출문제 (cst00309) artlSeq=5231525 file=raw/1008/5231525/2025년 제36회 감정평가사 1차시험 2교시 문제지 원본.pdf
%% conversion: kordoc 4.24.51 -> tools/gen1.py (auto); review=proofread-sonnet; flags=none
\begin{question}{8}
국토의 계획 및 이용에 관한 법령상 지구단위계획에 관한 설명으로 옳은 것은? (단, 조례는 고려하지 않음)
\begin{choices}
  \item 개발제한구역에서 해제되는 구역 중 계획적인 관리가 필요한 지역은 지구단위계획구역으로 지정할 수 없다.
  \item 지구단위계획구역의 지정목적이 한옥마을을 보존하고자 하는 경우에는 지구단위계획으로 ｢주차장법｣에 따른 주차장 설치기준을 150퍼센트까지 강화하여 적용한다.
  \item 도시지역에 개발진흥지구를 지정하고 당해 지구를 지구단위계획구역으로 지정한 경우에는 지구단위계획으로 ｢건축법｣ 제60조에 따라 제한된 건축물높이의 120퍼센트 이내에서 높이제한을 완화하여 적용할 수 있다.
  \item 주민이 입안을 제안한 지구단위계획에 관한 도시ㆍ군관리계획결정이 효력을 잃은 경우 해당 지구단위계획구역 지정 당시의 도시ㆍ군관리계획도 효력을 잃은 것으로 본다.
  \item 지구단위계획구역과 지구단위계획은 도시ㆍ군기본계획으로 결정한다.
\end{choices}
\end{question}
\answer{3}
